% ============================================================================
% DRAFT: t1_theorem_section.tex   (W2 deliverable, T1 landing draft)
% Date: 2026-08-04 / package: 2026-08-05
% Source: paper/current-usenix/t1_refinement_monotonicity_theorem_2026-08-04.md
% Status: INDEPENDENT DRAFT. Do NOT \input into main.tex until the PM
%         approves placement and the claim-map diff review has run.
%         This file does not modify any existing section.
%
% ----------------------------------------------------------------------------
% 落位建议 (placement suggestion)
%   File    : sections/security_analysis.tex
%   Position: new \subsection between "Counterfactual Witnesses and Bounded
%             Adequacy" (ends: "...surface and policy-irrelevant fields should
%             remain absent.") and "Conditional Effect Confinement".
%   Rationale:
%     - depends on the field-necessity criterion and collision-complete /
%       finite-adequacy machinery defined in that subsection;
%     - upgrades the appendix "Field necessity" paragraph from a single-point
%       argument to a monotone search statement (improvement_whitepaper T1);
%     - precedes "Conditional Effect Confinement" and does NOT alter the
%       O1--O5 premises of thm:single-commit (only cites O5 fail-closed).
%
% ----------------------------------------------------------------------------
% Dependencies (all existing labels; no new external citations needed)
%   used   : def:authorization-equivalence, prop:authorization-quotient,
%            prop:evidence-refinement, prop:ambiguous-cell,
%            def:collision-complete, thm:finite-adequacy;
%            method.tex §3 split/bind/trigger/expansion semantics and
%            "conservative registration or human review" fail-closed clause.
%   defined: sec:registration-convergence, def:refinement-lattice,
%            prop:refinement-operations, thm:refinement-order,
%            thm:refinement-termination, thm:refinement-collision-monotone,
%            cor:refinement-monotone-search.
%   companion edit (separate step, NOT included here): appendix/formal_proofs.tex
%            gains four proof paragraphs — skeletons are given in the trailing
%            comment block of this file.
%
% ----------------------------------------------------------------------------
% Number policy (数字策略)
%   WRITTEN numbers are frozen-domain results only (56-call domain + E2 grid):
%     trajectory SP 124->60->28->12->4->0; LB 37->33->25->9->4->0;
%     211 comparable lattice edges, 0 violations; length 5 <= |D|-1 = 55;
%     single-role deletion 16/16/16/8/4; common-contract baseline 118.
%   This subsection references NO v17 and NO V0-V3 numbers by construction.
% Frozen data sources (read-only):
%   experiments/security-analysis-ablation-and-overhead/results/
%     refinement-monotonicity-check/refinement-monotonicity-report.{json,md}
%   experiments/security-analysis-ablation-and-overhead/results/
%     policy-family-sensitivity/policy-family-sensitivity-report.md (cross-check)
%   NOTE: Table ref below (tab:policy-family-sensitivity) resolves only after
%   the companion draft e2_policy_family_table.tex is landed.
% ============================================================================

\subsection{Counterfactual Registration Converges}
\label{sec:registration-convergence}

The field-necessity criterion above is stated one field at a time. Registration,
however, searches a space of representations: the counterexample-guided loop of
Section~\ref{sec:security-analysis}'s method starts from a coarse candidate and
applies refinement operations until no authorization-separating collision
remains. We now show that, on a frozen finite domain and within the restricted
contract language, this search is a monotone descent with a bounded runtime.

\begin{definition}[Authorization partition and refinement order]
\label{def:refinement-lattice}
On a finite instance domain $D\subseteq\mathcal{U}$, let
\[
  Q(D)=\bigl\{[u]_{\equiv_{\mathcal{B}}}:u\in D\bigr\}
\]
be the authorization partition induced by
$\equiv_{\mathcal{B}}$ (Definition~\ref{def:authorization-equivalence}), and
let
\[
  P_\rho(D)=\bigl\{\rho^{-1}(z)\cap D:z\in\mathcal{Z}\bigr\}\setminus\{\emptyset\}
\]
be the partition induced by a representation $\rho$. Say that $\rho'$ refines
$\rho$ on $D$, written $\rho'\preceq\rho$, when
$\rho'(u)=\rho'(v)\Rightarrow\rho(u)=\rho(v)$ for all $u,v\in D$; equivalently,
every $P_{\rho'}$-cell is contained in a $P_\rho$-cell. The partitions of a
finite $D$ ordered by $\preceq$ form a finite lattice; every strict refinement
increases the number of cells by at least one, so every strict refinement
chain has length at most $|D|-1$.
\end{definition}

By Proposition~\ref{prop:authorization-quotient}, $Q(D)$ is the coarsest
authorization-sufficient representation, so registration searches for a
representable approximation to $Q(D)$ inside the sublattice reachable by the
restricted contract language. Let $\mathrm{SP}(\rho)$ be the set of separating
pairs---pairs $(u,v)$ with $\rho(u)=\rho(v)$ and $I_B(u)\neq I_B(v)$ for some
$B\in\mathcal{B}$---and write $|\mathrm{SP}(\rho)|$ for its cardinality. Let
$\mathrm{LB}(\rho)$ be the ambiguous-cell lower bound of
Proposition~\ref{prop:ambiguous-cell} applied to $P_\rho(D)$.

Registration refines through four operation classes: a \emph{split} separates
an effect label or attribute range, a \emph{bind} instantiates a symbolic or
wildcard attribute to a typed value, a \emph{trigger} records whether an
effect is conditioned on an argument or state predicate, and an
\emph{expansion} emits one occurrence per affected object or principal.

\begin{proposition}[Registered refinement operations refine]
\label{prop:refinement-operations}
Let $\rho'$ be obtained from $\rho$ by one admissible application of a split,
bind, trigger, or expansion refinement on the frozen domain $D$. If the
operation does not expand the underlying bound family $\mathcal{B}$ or the
authority view, then $\rho'\preceq\rho$.
\end{proposition}

\begin{proof}[Proof skeleton]
Verify $\rho'(u)=\rho'(v)\Rightarrow\rho(u)=\rho(v)$ by operation:
\begin{enumerate}
  \item split: $\rho'$ re-partitions each $\rho$-cell by a new attribute
        range; two contexts equal under $\rho'$ lie in the same sub-cell of a
        single $\rho$-cell;
  \item bind: binding injects the information that an attribute takes a
        specific value; $\rho'$-equality implies the same coarser binding
        state under $\rho$;
  \item trigger: the trigger predicate adds one boolean distinction;
        $\rho'$-equality implies equal trigger status, hence $\rho$-equality;
  \item expansion: expansion counts are part of the representation, and
        already-expanded objects are not re-expanded; $\rho'$-equality implies
        equal expansion counts, hence equality under the unexpanded
        representation. \qedhere
\end{enumerate}
\end{proof}

The legality condition is essential: free-form observations or model
assertions that expand the underlying bound are not refinements in this sense
(the restriction stated after Proposition~\ref{prop:evidence-refinement}), and
the results below do not apply to such steps.

\begin{theorem}[Registered refinement preserves partition order]
\label{thm:refinement-order}
Let $D\subseteq\mathcal{U}$ be a finite instance domain and let
$\rho_0,\rho_1,\ldots$ be a refinement trajectory generated by the restricted
contract language on $D$ in the sense of
Proposition~\ref{prop:refinement-operations}. Then for every step $i$,
\[
  \rho_{i+1}\preceq\rho_i,
\]
and consequently the separating-pair set is monotone non-increasing:
\[
  \mathrm{SP}(\rho_{i+1})\subseteq\mathrm{SP}(\rho_i),
  \qquad
  |\mathrm{SP}(\rho_{i+1})|\le|\mathrm{SP}(\rho_i)|.
\]
\end{theorem}

\begin{proof}[Proof skeleton]
The first claim is Proposition~\ref{prop:refinement-operations} applied step
by step. For the second, let $(u,v)\in\mathrm{SP}(\rho_{i+1})$: then
$\rho_{i+1}(u)=\rho_{i+1}(v)$, so $\rho_i(u)=\rho_i(v)$ by
$\rho_{i+1}\preceq\rho_i$, while the ideal decisions under $\mathcal{B}$ are
unchanged by representation-side operations; hence
$(u,v)\in\mathrm{SP}(\rho_i)$. This is the effect-side instance of the cell
nesting argument of Proposition~\ref{prop:evidence-refinement}.
\end{proof}

\begin{theorem}[Bounded termination of registered refinement]
\label{thm:refinement-termination}
Let $D$ be a finite instance domain with $|D|=n$, and let
$\rho_0,\rho_1,\ldots$ be a registered refinement trajectory that stops
exactly when it reaches a representation with no authorization-separating
collision or when no admissible refinement applies. Then the trajectory
terminates after at most $n-1$ strict refinements, in one of two states:
\begin{enumerate}
  \item a fixed point $\rho^\ast$ with $|\mathrm{SP}(\rho^\ast)|=0$, which is
        authorization-sufficient on $D$ whenever the validating suite is
        collision-complete (Definition~\ref{def:collision-complete},
        Theorem~\ref{thm:finite-adequacy}); or
  \item fail-closed routing to conservative registration or human review, the
        registration-side instance of obligation O5.
\end{enumerate}
\end{theorem}

\begin{proof}[Proof skeleton]
Every strict refinement increases the number of $P_\rho$-cells by at least one
(Definition~\ref{def:refinement-lattice}), and the number of cells is bounded
by $n$; hence at most $n-1$ strict refinements occur, and the trajectory
cannot cycle. If it stops with no remaining collision, the fixed point
satisfies $|\mathrm{SP}(\rho^\ast)|=0$, and a collision-complete suite then
certifies sufficiency on $D$ by Theorem~\ref{thm:finite-adequacy}. If it stops
because no admissible refinement eliminates a collision, the unresolved
mediation gap routes the tool to conservative registration or human review,
per the registration procedure; no $\Allow$-producing registration results.
\end{proof}

\begin{theorem}[Collision measures are monotone along refinement]
\label{thm:refinement-collision-monotone}
Along every trajectory of Theorem~\ref{thm:refinement-order},
\[
  |\mathrm{SP}(\rho_{i+1})|\le|\mathrm{SP}(\rho_i)|,
  \qquad
  \mathrm{LB}(\rho_{i+1})\le\mathrm{LB}(\rho_i).
\]
In particular, the number of authorized instances that any sound monitor must
withhold (Proposition~\ref{prop:ambiguous-cell}) never increases as
registration refines.
\end{theorem}

\begin{proof}[Proof skeleton]
The first inequality is Theorem~\ref{thm:refinement-order}. For the second,
every $\rho_{i+1}$-cell is contained in a $\rho_i$-cell; any authorized
instance counted by the $\rho_{i+1}$ lower bound lies in a mixed
$\rho_{i+1}$-cell whose unauthorized member is also contained in the enclosing
$\rho_i$-cell, so it was already counted by the $\rho_i$ lower bound. This
reuses the argument of Proposition~\ref{prop:evidence-refinement}.
\end{proof}

\begin{corollary}[Monotone minimal-sufficiency search]
\label{cor:refinement-monotone-search}
$\Phi(\rho)=(|\mathrm{SP}(\rho)|,\mathrm{LB}(\rho))$ is a non-increasing
potential along every registered refinement trajectory. A field or qualifier
is operationally necessary on $(D,\mathcal{B})$ if and only if erasing it
strictly increases $\Phi$; hence registration searching for a minimal
sufficient contract performs a monotone descent in $\Phi$ on the finite
partition lattice.
\end{corollary}

The corollary restates the appendix field-necessity criterion as a property of
the whole search path: a retained field is one whose deletion raises the
potential, and a field whose deletion leaves $\Phi$ unchanged is redundant on
$(D,\mathcal{B})$---the same criterion applied by the qualifier-deletion
sensitivity grid below.

\paragraph{Instantiation on the frozen 56-call domain.}
On the 56-call frozen domain, the refinement trajectory obtained by adding the
five typed qualifier roles in registration order is monotone non-increasing in
separating pairs, $124\to60\to28\to12\to4\to0$, with ambiguous-cell lower
bound $37\to33\to25\to9\to4\to0$; across all 211 comparable edges of the
five-role subset lattice, refinement order preservation and monotonicity of
both measures hold with zero violations, and the trajectory reaches the
zero-collision fixed point after 5 strict refinements, within the
$|D|-1=55$ bound of Theorem~\ref{thm:refinement-termination}. The
separating-pair decision procedure is exactly the one used by the
policy-family sensitivity grid (Table~\ref{tab:policy-family-sensitivity}):
single-role deletion yields $16/16/16/8/4$ separating pairs for
date/subject/recurrence/payload/visibility, matching the power-set rows, and
the check independently reproduces the frozen common-contract baseline of 118
separating pairs. One representation-semantics note: the lattice's fully
erased endpoint applies all five deletion operations and therefore starts at
124 separating pairs, whereas the common-contract baseline strips only the
entire qualifier key (118 pairs); these are two distinct coarse
representations, and both converge monotonically to the same typed endpoint.

\paragraph{Scope.}
Three qualifications are retained. First, legality: a step that expands the
underlying bound or authority view is not a refinement, and
Theorems~\ref{thm:refinement-order}--\ref{thm:refinement-collision-monotone}
do not apply to it. Second, relativity: all statements are relative to the
frozen domain $D$, the declared family $\mathcal{B}$, and the restricted
refinement language; convergence over an open tool domain or an unrestricted
synthesis language is not claimed. Third, the trajectory check above is
instantiation evidence on one frozen domain and one authority family, using
the E2 sensitivity data source; termination in general follows from finiteness
of the partition lattice, not from that single trajectory.

% ============================================================================
% COMPANION EDIT (appendix/formal_proofs.tex) — proof skeletons to be landed
% together with this subsection (T1 md §5.2, steps S2/S3/S4/S5):
%
% \paragraph{Proof of Proposition~\ref{prop:refinement-operations}.}
%   Four-case structural verification (split/bind/trigger/expansion) as in the
%   inline proof skeleton above; state the legality condition explicitly.
%
% \paragraph{Proof of Theorem~\ref{thm:refinement-order}.}
%   Cell nesting -> SP set inclusion; cite the
%   Proposition~\ref{prop:evidence-refinement} cell argument.
%
% \paragraph{Proof of Theorem~\ref{thm:refinement-termination}.}
%   Potential argument (strict refinement adds >=1 cell, cell count <= n),
%   fixed-point / fail-closed dichotomy, Theorem~\ref{thm:finite-adequacy}
%   for the sufficiency branch, O5 for the fail-closed branch.
%
% \paragraph{Proof of Theorem~\ref{thm:refinement-collision-monotone}.}
%   LB monotonicity via the evidence-refinement cell-containment argument.
%
% NON-COMMITMENT LIST (must survive porting, T1 md §6):
%   no open-tool-domain convergence; no sufficiency beyond collision-complete
%   suites; no refinement outside the restricted language; no monotonicity
%   when bounds are mutable.
% ============================================================================
